\documentclass[10pt]{beamer}

\usepackage[T1]{fontenc}

\usetheme{moloch}

\setbeamertemplate{page number in head/foot}[appendixframenumber]
\setbeamertemplate{section in toc}[sections numbered]

\usepackage{booktabs}
\usepackage[scale=2]{ccicons}

\usepackage[semibold,light]{FiraSans}
\usepackage{FiraMono}

\usepackage{xspace}

\newcommand{\themename}{\textbf{moloch}\xspace}

\title{Moloch}
\subtitle{A Minimal Beamer Theme}
\date{\today}
\author{The Author}
\institute{Some Institution, Some University}
\titlegraphic{\hfill\includegraphics[width=3.5cm]{moloch-logo.pdf}}

\begin{document}

\maketitle

\begin{frame}
\frametitle{Table of Contents}
  \tableofcontents[hideallsubsections]
\end{frame}

\section{Introduction}

\begin{frame}[fragile]
\frametitle{Moloch}

  The \themename theme is a Beamer theme with minimal visual noise. It is a fork of the
  \href{https://github.com/matze/mtheme}{metropolis theme} by Matthias Vogelgesang, which in
  turn was inspired by by the \href{https://github.com/hsrmbeamertheme/hsrmbeamertheme}{hsrm
    theme} by Benjamin Weiss. \medskip

  Enable the theme by calling
  \begin{verbatim}\documentclass{beamer}
\usetheme{moloch}\end{verbatim}
\end{frame}
\begin{frame}[fragile]
\frametitle{Sections}
  Sections group slides of the same topic by introducing a section page between them. A progress bar is shown which indicates how far along in the presentation you are.

  \begin{verbatim}\section{Title Formats}\end{verbatim}
\end{frame}

\section{Title Formats}

\begin{frame}[fragile]
\frametitle{Moloch Title Formats}
  \themename supports four different title formats:
  \begin{itemize}
    \item \texttt{regular} (Regular)
    \item \texttt{smallcaps} (\textsc{Small Caps})
    \item \texttt{allsmallcaps} (\textsc{all small caps})
    \item \texttt{allcaps} (ALL CAPS)
  \end{itemize}

  They can either be set globally for every frame or used locally just for the current frame
  and onwards by using
  \begin{verbatim}\molochset{titleformat frame=<option>}\end{verbatim}
\end{frame}

{
\molochset{titleformat frame=smallcaps}
\begin{frame}
\frametitle{Small Caps}

  \begin{alertblock}{Potential Problems}
    Be aware that not every font supports small caps. If you use the Computer (or Latin) Modern Sans Serif font, for instance, text in small caps will just be typeset in a normal font.
  \end{alertblock}
\end{frame}
}

{
\molochset{titleformat frame=allsmallcaps}
\begin{frame}
\frametitle{All Small Caps}
  This frame uses the \texttt{allsmallcaps} title format.

  \begin{alertblock}{Potential Problems}
    This title format also uses small caps, so you face the same problems as with the \texttt{smallcaps} title format.

    \medskip

    In addition, note that numbers and math will be unaffected by this setting and will match
    poorly with text as a result. For that reason, we suggest you make sure titles are all
    plain text if you use this option.

    \medskip

    Finally, observe that the height of the frame title box will be modified to account for the
    decreased text height since there are no capital letters. This too means that numbers and
    math will not work well.
  \end{alertblock}
\end{frame}
}

{
\molochset{titleformat frame=allcaps}
\begin{frame}
\frametitle{All Caps}
  This frame uses the \texttt{allcaps} title format.

  \begin{alertblock}{Potential Problems}
    This title format is not as problematic as the \texttt{allsmallcaps} format, but basically suffers from the same deficiencies. So please have a look at the documentation if you want to use it.
  \end{alertblock}
\end{frame}
}

\section{Elements}

\begin{frame}[fragile]
\frametitle{Typography}
  \begin{verbatim}The theme provides sensible defaults to
\emph{emphasize} text, \alert{accent} parts or show
\textbf{bold} results.\end{verbatim}
  \begin{center}
    becomes
  \end{center}
  The theme provides sensible defaults to \emph{emphasize} text,
  \alert{accent} parts or show \textbf{bold} results.
\end{frame}

\begin{frame}
\frametitle{Font Features Test}
  \begin{itemize}
    \item Regular
    \item \textit{Italic}
    \item \textbf{Bold}
    \item \textbf{\textit{Bold Italic}}
    \item \texttt{Monospace}
    \item \texttt{\textit{Monospace Italic}}
    \item \texttt{\textbf{Monospace Bold}}
    \item \texttt{\textbf{\textit{Monospace Bold Italic}}}
    \item \textsc{Small Caps}
    \item \textbf{\textsc{Bold Small Caps}}
  \end{itemize}
\end{frame}

\begin{frame}
\frametitle{Lists}
  \begin{columns}[T,onlytextwidth]
    \column{0.3\textwidth}
    Items
    \begin{itemize}
      \item Milk \item Eggs \item Potatoes
    \end{itemize}

    \column{0.33\textwidth}
    Enumerations
    \begin{enumerate}
      \item First, \item Second and \item Last.
    \end{enumerate}

    \column{0.33\textwidth}
    Descriptions
    \begin{description}[PowerPoint]
      \item[PowerPoint] Meeh. \item[Beamer] Yeeeha.
    \end{description}
  \end{columns}
\end{frame}
\begin{frame}
\frametitle{Animation}
  \begin{itemize}[<+- | alert@+>]
    \item \alert<4>{This is\only<4>{ really} important}
    \item Now this
    \item And now this
  \end{itemize}
\end{frame}
\begin{frame}
\frametitle{Figures}
  \begin{figure}
    \newcounter{density}
    \setcounter{density}{20}
    \begin{tikzpicture}
      \def\couleur{alerted text.fg}
      \path[coordinate] (0,0)  coordinate(A)
      ++( 90:5cm) coordinate(B)
      ++(0:5cm) coordinate(C)
      ++(-90:5cm) coordinate(D);
      \draw[fill=\couleur!\thedensity] (A) -- (B) -- (C) --(D) -- cycle;
      \foreach \x in {1,...,40}{%
          \pgfmathsetcounter{density}{\thedensity+20}
          \setcounter{density}{\thedensity}
          \path[coordinate] coordinate(X) at (A){};
          \path[coordinate] (A) -- (B) coordinate[pos=.10](A)
          -- (C) coordinate[pos=.10](B)
          -- (D) coordinate[pos=.10](C)
          -- (X) coordinate[pos=.10](D);
          \draw[fill=\couleur!\thedensity] (A)--(B)--(C)-- (D) -- cycle;
        }
    \end{tikzpicture}
    \caption{Rotated square from
      \href{http://www.texample.net/tikz/examples/rotated-polygons/}{texample.net}.}
  \end{figure}
\end{frame}

\begin{frame}
\frametitle{Tables}
  \framesubtitle{Largest cities in the world}
  \begin{table}
    \caption{Largest cities in the world (source: Wikipedia)}
    \begin{tabular}{@{} lr @{}}
      \toprule
      City        & Population \\
      \midrule
      Mexico City & 20,116,842 \\
      Shanghai    & 19,210,000 \\
      Peking      & 15,796,450 \\
      Istanbul    & 14,160,467 \\
      \bottomrule
    \end{tabular}
  \end{table}
\end{frame}
\begin{frame}[fragile]
\frametitle{Blocks}
  Three different block environments are pre-defined: \verb|block|, \verb|alertblock|,
  and \verb|exampleblock|. They can be styled using the option
  \verb|block|, which takes values \verb|transparent| (default, left column) and \verb|fill| (right column).

  \begin{columns}[T]
    \begin{column}{0.45\textwidth}
      \begin{block}{Default}
        Block content.
      \end{block}

      \begin{alertblock}{Alert}
        Block content.
      \end{alertblock}

      \begin{exampleblock}{Example}
        Block content.
      \end{exampleblock}
    \end{column}
    \begin{column}{0.45\textwidth}
      {
        \molochset{block=fill}

        \begin{block}{Default}
          Block content.
        \end{block}

        \begin{alertblock}{Alert}
          Block content.
        \end{alertblock}

        \begin{exampleblock}{Example}
          Block content.
        \end{exampleblock}
      }
    \end{column}
  \end{columns}
\end{frame}

\begin{frame}[c]
  \frametitle{Progress Bars}

  A prominent feature of \themename is the use of progress bars to
  visualize the progress of the presentation. \medskip

  By default bars are shown on section pages, but this can be toggled
  through \texttt{sectionpage=progressbar|simple|none} and also enabled for 
  subsection pages through \texttt{subsectionpage=progressbar}.\medskip

  
  Finally, you can enable progress bars on frames by setting 
  \texttt{progressbar=none|head|frametitle|foot} (depending on where
  you'd like the bar to be placed).\medskip
\end{frame}

\begin{frame}[c,fragile]
  \frametitle{Progress Bar Customization}
  You can customize the width of the bars using
  \texttt{progressbar linewidth=<length>}.\medskip

  On the following pages, we have set
\begin{verbatim}\molochset{
  progressbar=frametitle,
  progressbar linewidth=1pt
}
\end{verbatim}
  to enable progress bars below the frame title and set its width to 1pt.
\end{frame}

\molochset{progressbar=frametitle,progressbar linewidth=1pt}

\begin{frame}[c,fragile]
  \frametitle{The Title Page}

  There is only very little customization available of the title page at the
  time of writing.\medskip

  But you can at least modify the line width of the title separator
  by setting \texttt{titleseparator linewidth=<length>}
\end{frame}

\begin{frame}
\frametitle{Math}
  \begin{equation*}
    e = \lim_{n\to \infty} \left(1 + \frac{1}{n}\right)^n
  \end{equation*}
\end{frame}

\begin{frame}
\frametitle{Quotes}
  \begin{quote}
    Verily, I say unto you, the days spoken of in the Apocalypse are nigh!
  \end{quote}
\end{frame}

{%
\setbeamertemplate{frame footer}{My custom footer}
\begin{frame}[fragile]
\frametitle{Frame Footer}
  \themename defines a custom beamer template to add a text to the footer. It can be set via
  \begin{verbatim}\setbeamertemplate{frame footer}{My custom footer}\end{verbatim}
\end{frame}
}

\begin{frame}[fragile]
\frametitle{References}
  Here are some references~\cite{Knuth92,ConcreteMath,Simpson,Er01,greenwade93} to showcase \verb+[allowframebreaks]+.
\end{frame}

\section{Conclusion}

\begin{frame}
\frametitle{Summary}

  Get the source of this theme and the demo presentation from
  \begin{center}
    \url{github.com/jolars/moloch}
  \end{center}

  The theme is licensed under the
  \href{http://creativecommons.org/licenses/by-sa/4.0/}{Creative Commons
    Attribution-ShareAlike 4.0 International License}.

  \begin{center}
    \ccbysa
  \end{center}

\end{frame}

\begin{frame}[standout]
  Questions?
\end{frame}

\appendix

\begin{frame}[allowframebreaks]{References}
  \bibliography{demo}
  \bibliographystyle{abbrv}
\end{frame}

\end{document}
